\begin{abstract}
%{\small
% Introduction, background
Real-world software applications must constantly evolve to remain relevant. This evolution occurs when developing new applications or adapting existing ones to meet new requirements, make corrections, or incorporate future functionality.
Traditional methods of software quality control involve software quality models and continuous code inspection tools. These measures focus on directly assessing the quality of the software.
However, there is a strong correlation and causation between the quality of the development process and the resulting software product.
Therefore, improving the development process indirectly improves the software product, too.
To achieve this, effective learning from past processes is necessary, often embraced through post-mortem organizational learning.
While qualitative evaluation of large artifacts is common, smaller quantitative changes captured by application lifecycle management are often overlooked.
In addition to software metrics, these smaller changes can reveal complex phenomena related to project culture and management.
Leveraging these changes can help detect and address such complex issues.


% Problem?
% Method, how did we try to solve the problem?
% Results/Analysis, what was the outcome?
% Discussion, what do the results mean?
Software evolution was previously measured by the size of changes, but the lack of consensus on a reliable and versatile quantification method prevents its use as a dependable metric.
Different size classifications fail to reliably describe the nature of evolution. While application lifecycle management data is rich, identifying which artifacts can model detrimental managerial practices remains uncertain.
Approaches such as simulation modeling, discrete events simulation, or Bayesian networks have only limited ability to exploit continuous-time process models of such phenomena.
Even worse, the accessibility and mechanistic insight into such gray- or black-box models are typically very low.
%Simulation modeling and Bayesian networks have limited ability to exploit continuous-time process models, with low accessibility and mechanistic insight.
To address these challenges, we suggest leveraging objectively captured digital artifacts from application lifecycle management, combined with qualitative analysis, for efficient organizational learning.
A new language-independent metric is proposed to robustly capture the size of changes, significantly improving the accuracy of change nature determination.
The classified changes are then used to explore, visualize, and suggest maintenance activities, enabling solid prediction of malpractice presence and -severity, even with limited data.
Finally, parts of the automatic quantitative analysis are made accessible, potentially replacing expert-based qualitative analysis in parts.
%
%
\bigskip
\vfill
\noindent
\textbf{Keywords:}
Software Size, Software Metrics, Commit Classification, Maintenance Activities, Software Quality, Process Quality, Project Management, Organizational Learning, Machine Learning, Visualization, Optimization

\end{abstract}
\cleardoublepage
